\chapter{Anexos}

\section{Prompt usado para entrevista simulada}

\begin{summarybox}
Actua como Carlos Pérez, trabajador con mas de diez años de experiencia en Alestis Puerto Real. Lunabotics necesita preparar una oferta tecnica y comercial de robotica movil para mejorar procesos relacionados con el HTP del A320 y la seccion 19.1. Responde como experto de planta: explica que ocurre en esa zona, que materiales o herramientas se mueven, que problemas recurrentes aparecen, si hay retrasos por componentes o utiles no localizados, que papel tiene FOD y que restricciones tendria un robot movil. Si la conversacion se desvia hacia temas personales, Cadiz o cultura local, redirige de forma natural hacia el proceso productivo. No propongas que el robot fabrique directamente el HTP; enfocate en transporte interno, kitting, retorno de utiles, trazabilidad RFID/QR, inspeccion preventiva y datos utiles para una oferta comercial.
\end{summarybox}

\section{Capturas de entrevista}

\begin{itemize}
  \item Placeholder 1: captura de inicio de conversacion con Carlos Pérez.
  \item Placeholder 2: captura de respuestas sobre HTP, seccion 19.1 y flujo logistico.
  \item Placeholder 3: captura de respuestas sobre FOD, trazabilidad y automatizacion.
\end{itemize}

\section{Matriz de cumplimiento contra la guia oficial}

\begin{longtable}{P{0.32\textwidth}P{0.46\textwidth}P{0.14\textwidth}}
\toprule
\textbf{Exigencia de la guia} & \textbf{Evidencia en el documento} & \textbf{Estado}\\
\midrule
\endhead
Actividad 1 de Sistemas Roboticos Moviles, con entrega el 26 de mayo de 2026 antes de las 23:59. & Portada, metadatos del informe y metodologia. & Cumplido\\
Diseño teorico de un sistema robotico para una aplicacion industrial concreta. & Propuesta AMR-FOD-Kitting para Alestis Puerto Real, HTP A320 y seccion 19.1. & Cumplido\\
Parte 1: entrevista con un experto. & Capitulo de entrevista con ficha, acta tecnica, conversacion simulada, hallazgos e implicaciones. & Cumplido\\
Parte 2: oferta tecnica basada en robotica movil. & Identificacion del problema, requisitos, alternativas, robot seleccionado, sensores, arquitectura, operacion y validacion en CoppeliaSim. & Cumplido\\
Parte 3: oferta comercial de la misma solucion. & Presupuesto por partidas, escenarios 1/7/17 robots, TCO, ROI y condiciones comerciales academicas. & Cumplido\\
Introduccion a Lunabotics. & Introduccion y presentacion de la empresa ficticia en informe y slides. & Cumplido\\
Identificacion del problema en Alestis Puerto Real. & Capitulos de entrevista, problema, requisitos y matriz de trazabilidad. & Cumplido\\
Propuesta de solucion con robot movil disponible en CoppeliaSim. & AMR omnidireccional tipo KUKA YouBot/Omnirob o equivalente. & Cumplido\\
Viabilidad tecnica y comercial, incluyendo beneficios y retorno. & Capitulos de viabilidad, riesgos, presupuesto y ROI. & Cumplido\\
Datos clave: 40 HTP/mes, objetivo 80 HTP/mes o mas, salario 20.000--28.000 euros/año. & Metodologia, despliegue y analisis economico. & Cumplido\\
Añadir una diapositiva extra detallando presupuesto. & Diapositiva \textit{Presupuesto detallado} en presentacion PDF y PPTX. & Cumplido\\
Incluir tipo de robot y sensores adicionales. & Capitulo de propuesta tecnica, arquitectura, anexos y slides. & Cumplido\\
Especificar cada partida presupuestaria. & Tabla de presupuesto por partidas y CSV auxiliar. & Cumplido\\
\bottomrule
\end{longtable}

\section{Tabla completa de presupuesto}

La Tabla~\ref{tab:presupuesto} contiene la version completa del presupuesto por partidas para 1, 7 y 17 robots. Esta tabla se replica de forma resumida en la diapositiva extra de presupuesto de la presentacion.

\section{Diapositiva extra de presupuesto}

La presentacion incluye una diapositiva especifica titulada \textit{Presupuesto detallado}. Su objetivo es cumplir el requisito del enunciado de añadir una diapositiva extra con las partidas presupuestarias.

\section{Lista de sensores}

\begin{itemize}
  \item LiDAR 2D/3D para SLAM, obstaculos y seguridad.
  \item Camara RGB-D o industrial para inspeccion FOD.
  \item Lector RFID/QR para trazabilidad.
  \item IMU y encoders para odometria.
  \item Bumpers, escaner laser de seguridad y parada de emergencia.
  \item HMI/tablet y comunicacion WiFi industrial.
\end{itemize}

\section{Checklist de cumplimiento de la actividad}

\begin{longtable}{P{0.68\textwidth}P{0.22\textwidth}}
\toprule
\textbf{Criterio} & \textbf{Estado}\\
\midrule
\endhead
Portada con asignatura, actividad, profesor y fecha correcta. & Cumplido\\
Mencion de Alestis Puerto Real, HTP, seccion 19.1 y A320. & Cumplido\\
Entrevista simulada con formato dedicado y tabla de hallazgos. & Cumplido\\
Oferta tecnica clara basada en robotica movil. & Cumplido\\
Oferta comercial con presupuesto por partidas. & Cumplido\\
Escenarios de 1, 7 y 17 robots. & Cumplido\\
Robot seleccionado y sensores adicionales. & Cumplido\\
Uso de CoppeliaSim como marco de validacion. & Cumplido\\
Presentacion PDF y PPTX editable equivalente a PowerPoint. & Cumplido\\
Diapositiva extra de presupuesto. & Cumplido\\
Eliminacion de contenido ajeno al alcance de la Actividad 1. & Cumplido\\
\bottomrule
\end{longtable}

\section{Checklist por ponderacion}

\begin{longtable}{P{0.18\textwidth}P{0.50\textwidth}P{0.22\textwidth}}
\toprule
\textbf{Bloque} & \textbf{Evidencia incluida} & \textbf{Cobertura}\\
\midrule
\endhead
Entrevista 20\% & Ficha, acta tecnica, conversacion simulada con Carlos Pérez, hallazgos e implicaciones directas para la solucion. & Completa\\
Oferta tecnica 60\% & Problema, requisitos, alternativas, robot seleccionado, sensores, arquitectura, operacion, CoppeliaSim, escenarios, riesgos y KPIs. & Completa\\
Oferta comercial 20\% & Presupuesto por partidas para 1, 7 y 17 robots, TCO, ROI prudente y condiciones comerciales academicas. & Completa\\
\bottomrule
\end{longtable}

\section{Artefactos de entrega}

\begin{itemize}
  \item Informe PDF: \texttt{actividad1/main.pdf}.
  \item Presentacion PDF: \texttt{actividad1/slides/slides.pdf}.
  \item Presentacion editable PowerPoint: \texttt{actividad1/slides/Actividad1\_AMR\_FOD\_Kitting.pptx}.
  \item Tablas auxiliares: carpeta \texttt{actividad1/tables/}.
\end{itemize}
